% id: 40-02
% book: 베이직쎈 확률과 통계 · 03 확률의 뜻과 활용
% section: 유형 022 합사건, 곱사건, 여사건
% figure: none
% answer: 
% answer_source: 
% variant_level: 0 (원본 전사)
% 이 파일은 build.mjs 가 items.json 에서 만든다. 고칠 때는 items.json 을 고친다.
\smash{\raisebox{1.5mm}{\dmkichul{베이직쎈 확률과 통계 40쪽 02번}}}
각 면에 1부터 12까지의 자연수가 각각 하나씩 적힌 정십이면체를 던지는 시행에서 바닥에 닿는 면에 적힌 수가 4의 배수인 사건을 $A$, 8 이하의 수인 사건을 $B$라 할 때, 다음 중 옳지 \underline{않은} 것은?
\par\medskip\par\noindent\hangindent=1.4em\hangafter=1 {\small ①}\ $A=\{4{,}\mkern3mu\mathopen{}8{,}\mkern3mu\mathopen{}12\}$\par\noindent\hangindent=1.4em\hangafter=1 {\small ②}\ $A\cap B=\{4{,}\mkern3mu\mathopen{}8\}$\par\noindent\hangindent=1.4em\hangafter=1 {\small ③}\ $A\cup \comp{B}=\{4{,}\mkern3mu\mathopen{}8{,}\mkern3mu\mathopen{}9{,}\mkern3mu\mathopen{}10{,}\mkern3mu\mathopen{}11{,}\mkern3mu\mathopen{}12\}$\par\noindent\hangindent=1.4em\hangafter=1 {\small ④}\ $n(\comp{A})=9$\par\noindent\hangindent=1.4em\hangafter=1 {\small ⑤}\ $n(\comp{A}\cap B)=5$\par\medskip
